\section{Results}
\vspace{-3mm}
\label{sec: results}
With contextual snapshots as input, we evaluate different LLMs by simulating frozen decision points, where the agent has interacted with the environment under a specific task instruction and just received the latest observation. Thus, allowing for the reproducible assessment of the faithfulness of the generated current step action.

We benchmark twelve prominent instruction-tuned LLMs, spanning both open-source and proprietary families, varying in size and version, including the latest reasoning-oriented models. Deterministic decoding with temperature set to 0 is used throughout all evaluations to ensure reproducibility.

\newcommand{\resulttablestyle}[2]{\setlength{\tabcolsep}{#1}\renewcommand{\arraystretch}{#2}\centering\footnotesize}

\newcommand\scline[1]{\noalign{\global\savewidth\arrayrulewidth
  \global\arrayrulewidth .8pt}\cline{#1}\noalign{\global\arrayrulewidth\savewidth}}

% table 4
\begin{table*}
    \centering
    \resulttablestyle{3pt}{1.3}
    \resizebox{\textwidth}{!}{%
    \begin{tabular}{lcccccccccccccccc}
    \thickcline{1-17}
    \multirow{2}{*}{\large\textbf{Models}} 
    & \multicolumn{2}{c}{\cellcolor{cyan!5}\textbf{OSQ}} 
    & \multicolumn{2}{c}{\cellcolor{orange!10}\textbf{UET}} 
    & \multicolumn{2}{c}{\cellcolor{blue!10}\textbf{UGS}} 
    & \multicolumn{2}{c}{\cellcolor{pink!10}\textbf{ISI}} 
    & \multicolumn{2}{c}{\cellcolor{teal!10}\textbf{FIH-R}} 
    & \multicolumn{2}{c}{\cellcolor{violet!10}\textbf{FIH-E}} 
    & \multicolumn{2}{c}{\cellcolor{brown!15}\textbf{PUD}} 
    & \multicolumn{2}{c}{\cellcolor{gray!10}\textbf{Overall}} \\
    & US\,$\uparrow$ & HR\,$\downarrow$ 
    & US\,$\uparrow$ & HR\,$\downarrow$ 
    & US\,$\uparrow$ & HR\,$\downarrow$ 
    & US\,$\uparrow$ & HR\,$\downarrow$ 
    & US\,$\uparrow$ & HR\,$\downarrow$ 
    & US\,$\uparrow$ & HR\,$\downarrow$ 
    & US\,$\uparrow$ & HR\,$\downarrow$ 
    & US\,$\uparrow$ & HR\,$\downarrow$ \\


    \thickcline{1-17}
    \rowcolor{gray!20}\multicolumn{17}{l}{\textit{\textbf{Open-source Models}}} \\

        Qwen2.5-7B-Instruct & 0.409 & 0.496 & 0.385 & 0.551 & 0.225 & 0.557 & 0.393 & 0.493 & 0.595 & 0.351 & 0.533 & 0.377 & 0.995 & \textbf{0.000} & 0.488 & 0.416 \\
    Qwen2.5-32B-Instruct & 0.487 & 0.407 & 0.526 & 0.462 & 0.336 & 0.414 & 0.514 & 0.333 & 0.615 & 0.330 & \textbf{0.725} & \textbf{0.238} & 0.995 & \textbf{0.000} & 0.581 & 0.324 \\
    Qwen2.5-72B-Instruct & 0.544 & 0.371 & 0.577 & 0.372 & 0.264 & 0.500 & 0.518 & 0.362 & 0.632 & 0.312 & 0.642 & 0.325 & 0.995 & \textbf{0.000} & 0.579 & 0.334 \\
    Llama-3.1-70B-Instruct & 0.470 & 0.453 & 0.378 & 0.564 & 0.261 & 0.557 & 0.409 & 0.435 & 0.395 & 0.591 & 0.583 & 0.371 & 0.995 & \textbf{0.000} & 0.480 & 0.442 \\
    Llama-3.3-70B-Instruct & 0.420 & 0.496 & 0.397 & 0.564 & 0.186 & 0.650 & 0.438 & 0.391 & 0.531 & 0.438 & 0.583 & 0.377 & 0.980 & 0.010 & 0.488 & 0.433 \\
    \textcolor{gray}{DeepSeek-chat} & \textcolor{gray}{0.622} & \textcolor{gray}{0.266} & \textcolor{gray}{0.520} & \textcolor{gray}{0.459} & \textcolor{gray}{0.248} & \textcolor{gray}{0.600} & \textcolor{gray}{0.477} & \textcolor{gray}{0.415} & \textcolor{gray}{0.466} & \textcolor{gray}{0.417} & \textcolor{gray}{0.624} & \textcolor{gray}{0.338} & \textcolor{gray}{0.979} & \textcolor{gray}{0.021} & \textcolor{gray}{0.549} & \textcolor{gray}{0.362} \\
    \textcolor{gray}{DeepSeek-reasoner} & \textcolor{gray}{0.511} & \textcolor{gray}{0.428} & \textcolor{gray}{0.493} & \textcolor{gray}{0.365} & \textcolor{gray}{0.533} & \textcolor{gray}{0.348} & \textcolor{gray}{0.518} & \textcolor{gray}{0.053} & \textcolor{gray}{0.687} & \textcolor{gray}{0.213} & \textcolor{gray}{0.886} & \textcolor{gray}{0.070} & \textcolor{gray}{0.963} & \textcolor{gray}{0.021} & \textcolor{gray}{0.641} & \textcolor{gray}{0.257} \\

    \rowcolor{gray!20}\multicolumn{17}{l}{\textit{\textbf{Proprietary Models}}} \\

    Claude-3-5-sonnet-20240620 & 0.524 & 0.341 & 0.539 & 0.390 & 0.291 & 0.631 & 0.589 & 0.244 & \textbf{0.746} & \textbf{0.227} & 0.565 & 0.253 & 0.892 & 0.082 & \textbf{0.589} & 0.308 \\
    Claude-3.7-Sonnet & 0.498 & 0.371 & \textbf{0.686} & \textbf{0.205} & 0.336 & 0.457 & \textbf{0.605} & \textbf{0.203} & 0.624 & 0.325 & 0.662 & 0.185 & 0.840 & 0.150 & 0.585 & \textbf{0.290} \\
    Gemini-2.0-flash & 0.532 & 0.345 & 0.532 & 0.462 & 0.225 & 0.614 & 0.406 & 0.486 & 0.641 & 0.348 & 0.599 & 0.291 & \textbf{1.000} & \textbf{0.000} & 0.547 & 0.369 \\
    Gemini-2.5-flash & \textbf{0.580} & 0.328 & 0.558 & 0.269 & \textbf{0.353} & \textbf{0.412} & 0.504 & 0.341 & 0.630 & 0.333 & 0.599 & 0.325 & 0.974 & 0.021 & 0.586 & 0.308 \\
    GPT-4o-mini-2024-07-18 & 0.489 & 0.366 & 0.545 & 0.397 & 0.214 & 0.571 & 0.511 & 0.377 & 0.503 & 0.435 & 0.540 & 0.404 & 0.990 & 0.010 & 0.518 & 0.381 \\
    GPT-4o-2024-11-20 & 0.560 & \textbf{0.323} & 0.577 & 0.385 & 0.243 & 0.593 & 0.511 & 0.348 & 0.572 & 0.345 & 0.652 & 0.311 & 0.990 & \textbf{0.000} & 0.569 & 0.339 \\

    \bottomrule
    \end{tabular}
    }
    \caption{Utility Scores (US) and Hallucination Rates (HR) of \textit{Out of Scope Queries (USQ)}, \textit{Unexpected Environmental Transitions (UET)}, \textit{Unachievable Goal State (UGS)}, \textit{Ill-specified Instructions (ISI)}, \textit{Flawed Interaction History-repetitive (FIH-R)}, \textit{Flawed Interaction History-erroneous (FIH-E)}, \textit{Pop-up Distractions (PUD)}. Scores for \textit{DeepSeek-reasoner} and \textit{DeepSeek-chat} are shown in gray to indicate evaluation was incomplete due to context window constraints, as some snapshots contain as many as 10k tokens.}
    \label{tab:overall_results}
    \vspace{-3mm}
\end{table*}
 
\paragraph{Persistent challenge from hallucination.}
Despite differences in architecture and scale, all evaluated models exhibit relatively low utility scores (\textit{US} typically below 0.6) and substantial hallucination rates (\textit{HR} often above 0.3), highlighting persistent reliability issues under realistic, context-sensitive perturbations. Even top proprietary models, such as \textit{GPT-4o-2024-11-20} (\textit{US}=0.569, \textit{HR}=0.339) and \textit{Gemini-2.5-flash} (\textit{US}=0.586, \textit{HR}=0.308), frequently produce unfaithful actions, underscoring the ongoing severity of hallucination issues across model types.


\paragraph{Marginal lead from proprietary models.}
Although proprietary models generally outperform open-source counterparts, the performance gap remains modest. Notably, the open-source model \textit{Qwen2.5-32B-Instruct} (\textit{US}=0.581, \textit{HR}=0.324) achieves competitive scores with proprietary models like \textit{GPT-4o-2024-11-20}. This indicates that scaling and advanced instruction tuning alone do not fully resolve faithfulness issues, emphasizing the need for targeted alignment methods specific to interactive agent contexts. To confirm these performance gaps are robust, stability analysis on three representative models shows consistent hallucination and utility scores with low variance (Table~\ref{tab:stability-summary}).

\paragraph{Presumptive hallucination pattern.}
Among various hallucination patterns identified in our detailed analysis of agent behaviors under each risk setting (Appendix~\ref{app:risk_settings}), we highlight one recurring issue here: LLM agents frequently fabricate contextually unsupported information, reminiscent of chatbot-style generation. For instance, when a required button is not observed, the agent may assume its existence and attempt to click it. Similarly, in the absence of an interactive user response, the agent may hallucinate a reply and proceed accordingly. Such behaviors, rarely seen in human decision-making, likely stem from instruction tuning in open-domain dialogue settings, which encourages speculative completions---an inductive bias misaligned with the demands of agentic environments.

\paragraph{Pop-up distractions: minor susceptibility emerges in stronger models.}
Stronger models (e.g., \textit{Claude-3.5-Sonnet}) exhibit slightly higher rates (0.08) of interaction with pop-up distractions, while weaker models almost never do---suggesting that increased perceptual capacity may introduce mild susceptibility to irrelevant cues and lead to hallucination. Still, most models show surprising robustness to pop-ups in our setting. Unlike prior work in vision~\cite{ma2024caution, yang2025riosworld}, where agents are easily distracted, our text-based pop-ups are injected via accessibility trees, matching real OSWorld observations. This gap suggests that certain risk settings may be alleviated by optimizing the modality and structure of agent observations, as well as other aspects of agent framework design.